% ------------------------------------------------------------------
%  Results.  Declarative subsection headings; each states a finding.
%  Every number that comes from the student slides rather than from the
%  evaluation JSON is wrapped in \todo{} (see README, "Before submission").
% ------------------------------------------------------------------

\subsection*{\toolname separates evidence extraction, rule-based selection and literature-grounded rescue}

\toolname takes, for one patient, the standardized record of nine sources: admission diagnosis, underlying conditions, blood counts, other laboratory values, cultures, \filmarray results, galactomannan, imaging reports and the laboratory's \mngs report with per-organism read counts and no-template-control codes (Fig.~\ref{fig:overview}a). Each modality is read by its own language-model agent into a fixed JSON schema. The agents are single-purpose: the imaging agent sees only imaging, the culture agent only cultures, and none of them sees the answer labels. A deterministic summary then joins the per-modality outputs into organism-level support and a host-vulnerability tier.

Selection is deterministic (Fig.~\ref{fig:overview}b). Every organism on the \mngs report is ranked by its control-subtracted priority and read percentile, tiered by read burden (R0--R4, decade thresholds from 1 to 10,000 reads) and by dominance within the specimen (D3 when the top organism has at least ten times the reads of the runner-up), aligned with the target site, and joined with same-organism hospital evidence. These features map to an \mngs signal tier (M1 strong to M5 background) and, together with site alignment and host status, to one of five causative levels. Pre-specified guardrails demote respiratory \textit{Candida}, oral and upper-airway flora, water-borne Gram-negative environmental organisms, skin flora, Epstein--Barr virus and other low-actionability herpesviruses, weak mould signals and generic hospital labels. Level~1--2 organisms and level~3 organisms with a protected or supported rationale are picked, in a fixed order that no later stage may change. A language-model safety review then examines only the unpicked organisms and sorts them into high-priority, context-needed, low-specificity or omitted; the first two tiers, and only those, are carried forward.

For every carried candidate, four evidence modules are computed (Fig.~\ref{fig:overview}c). Module A retrieves up to ten PubMed articles for the organism at the target site and judges each for exact-species, human, site-matched causal evidence with a quoted supporting span; module A1 re-judges every article against the patient---site, syndrome, host, phenotype and timing---and returns strong match, partial match, mismatch or insufficient. Module B scores \mngs strength on an absolute axis (signal tier M1--M2 or read tier R3--R4) and a relative axis (dominance D2--D3, or rank~1 with read percentile $\ge 0.8$). Module C grades same-organism non-\mngs evidence from C0 (none) to C3 (exact organism in a sterile-site or tissue specimen), and module D applies a specimen-colonization gate. Missing or pending tests are treated as unknown, never as negative. The final rule \rfive (Fig.~\ref{fig:overview}d) keeps every upstream pick and rescues a candidate only if C3 holds, or if A1 (at least two strong matches, or at least seven strong plus partial) coincides with site alignment and at least one B axis. Only then is a rationale written, by a language model constrained to the frozen selection and the patient's evidence window, and audited for structure before delivery as a simple table and a traceable JSON, Markdown and XLSX package.

\subsection*{\toolname identifies causative pathogens more accurately than any single conventional read-out}

We evaluated \toolname on \todo{55} patients with suspected lower respiratory tract infection from Kaohsiung Medical University Hospital (\kmuh, $n=\todo{30}$), Taipei Veterans General Hospital (\tvgh) and Tri-Service General Hospital (\tsgh; together $n=\todo{25}$), drawn from a three-hospital \mngs registry of 105 specimens (Table~\ref{tab:cohort}; Supplementary Fig.~1). Patients were predominantly critically ill (71\% in intensive care at sampling), 56\% had hospital- or ventilator-associated pneumonia, 37\% had an active malignancy and 39\% were receiving systemic corticosteroids; 60\% of \mngs specimens were bronchoalveolar lavage fluid. The reference standard was the causative pathogen list adjudicated by the treating infectious-diseases physicians, with versioned label revisions (Methods). Because a patient may have more than one pathogen, we scored (patient, organism) pairs at exact species level with a frozen synonym table and report micro-averaged precision, recall and F1.

\toolname reached a precision of \todo{71.6\%}, a recall of \todo{91.4\%} and an F1 of \todo{80.3\%} (Fig.~\ref{fig:benchmark}a,b). The raw \mngs report, scored by taking every listed organism as a call, had the highest recall (\todo{93.1\%}) but a precision of \todo{8.0\%} (F1 \todo{14.8\%}): the report contains the pathogen almost every time, but for every causative organism it names it lists about \todo{eleven} that are not \todo{(1--0.08)/0.08; confirm FP count from the evaluation JSON}. The targeted tests behaved in the opposite way. \filmarray and galactomannan together had a precision of \todo{67.6\%} but a recall of \todo{39.7\%} (F1 \todo{50.0\%}), and culture reached \todo{48.1\%} precision and \todo{43.1\%} recall (F1 \todo{45.5\%}). \toolname therefore kept essentially all of the sensitivity of \mngs while raising its precision by \todo{ninefold}, and it exceeded the precision of the targeted tests without their blind spots for fungi, viruses and fastidious organisms. \todo{Add per-hospital metrics, exact-set-match rate and any-hit rate from the evaluation JSON; add 95\% confidence intervals by patient-level bootstrap.}

\subsection*{\toolname reduces the organism burden presented to the physician}

The clinical cost of \mngs is not only false attribution but the length of the list a physician must reconcile. In an earlier series of \todo{35} intensive-care patients from \kmuh evaluated with the first version of the pipeline, culture reported at least one organism in \todo{24\%} of patients and \filmarray in \todo{64\%}, whereas \mngs reported at least one organism in every patient, with a mean of \todo{5.76} organisms per patient (Fig.~\ref{fig:benchmark}c,d). \toolname retained a call in every patient but presented a mean of \todo{1.97} organisms, close to the \todo{one to two} pathogens per patient in the adjudicated reference. The reduction is not a fixed cut-off: it comes from the demotion of colonizers, contaminants and reactivation viruses that fail the site, dominance and support rules, and from the guardrails that hold respiratory \textit{Candida} and oral flora below the picking threshold unless independent evidence exists. \todo{Confirm the relation of this 35-patient series to the 55-patient cohort and recompute organisms per patient for the 55 from the R5 decisions (518 candidates $\rightarrow$ 123 selected, i.e. 9.4 $\rightarrow$ 2.2 per patient).}

\subsection*{Structured integration outperforms direct prompting of frontier language models}

To test whether the gain comes from the structure of the pipeline rather than from the language models inside it, we gave four frontier models (\gptsol, \gptluna, \gptterra and \gptastra) exactly the raw data that \toolname receives---the \mngs list with reads and control codes, cultures, \filmarray, galactomannan, molecular tests, imaging reports, blood counts, diagnosis and underlying conditions---for \todo{41} patients (\kmuh $n=\todo{30}$, \tvgh/\tsgh $n=\todo{11}$), de-identified, with a single expert-written prompt that explained the colonization and site rules in words, a JSON output schema and medium reasoning effort, and no retrieval, candidate pool or answer key (Methods; Supplementary Note~\todo{2}). Predictions were frozen before the reference labels were read.

All four models converged on the same operating point: precision between \todo{67.7\%} and \todo{70.5\%}, recall between \todo{56.6\%} and \todo{60.5\%}, and F1 between \todo{62.4\%} and \todo{63.9\%} (Fig.~\ref{fig:llm}). On the same \todo{41} patients, \toolname reached a recall of \todo{90.8\%} at a precision of \todo{67.7\%} (\todo{69} true positives, \todo{33} false positives, \todo{7} false negatives), for an F1 of \todo{77.5\%}, \todo{13.6--15.1} points above the best direct-prompting model. The difference is almost entirely recall: prompted with the same warnings about colonization that \toolname encodes as rules, the models declined to call organisms that the rules, the hospital evidence and the case-matched literature together supported, whereas their precision was no better than the pipeline's. \todo{Slide 3 lists the pipeline's precision on these 41 patients as 71.62\%; the frozen tally TP 69 / FP 33 gives 67.65\%, and the slide's recall and F1 match that tally. Resolve from reports/idfixed\_20260906 before submission.} The newest model (\gptastra) improved recall by only \todo{2.6} points over its predecessors, which suggests that the shortfall is not a matter of model capacity but of what a single prompt can be made to do with a long, noisy organism list.

\subsection*{Rescue by case-matched literature contributes a minority of selections}

Across the \todo{55} patients, the deterministic scorer considered \todo{518} organism candidates and picked \todo{114}; the \rfive rule rescued a further \todo{9}, giving \todo{123} selected organisms (\todo{2.2} per patient). Rescue therefore changed \todo{7\%} of the final list, and every rescued organism carries its rule path (C3 direct, or A1 threshold with site alignment and B axis) and the PubMed identifiers of the articles that met the case-fit thresholds. \todo{From the evaluation JSON: how many of the 9 rescued organisms were true positives; how many false negatives remained and why (not in report, below read tier, guardrail); how many false positives were upstream picks versus rescues.} The immutability of upstream picks means that rescue can only add recall; the precision of the final list is bounded by the deterministic stage.

\subsection*{Decisions are traceable to their evidence}

Two cases illustrate the output (Supplementary Note~\todo{3}; details anonymized). In a patient with new bilateral consolidation and pleural effusions, a rising C-reactive protein and worsening hypoxaemia, the \mngs report listed \textit{Klebsiella pneumoniae} with \todo{68,569} reads alongside \textit{Candida albicans}; \toolname selected \textit{K.~pneumoniae} (dominant, aligned, supported by imaging and inflammatory markers) and excluded \textit{C.~albicans} as respiratory colonization with low read counts and repeated low-burden culture. In a second patient, \textit{Candida tropicalis} was selected despite the general demotion of respiratory yeasts, because multiple blood-culture sets were positive for the same species and the bronchoalveolar \mngs reads were high, which satisfies the C3 direct-evidence path. \todo{Both vignettes come from the January 2026 prototype (talk slide 15); regenerate them from the AETIA R5 rationale outputs for the same patients.} In each case the delivered record names the rule that fired, the modules that supported or opposed each organism, and the SHA-256 hash of every input file, so that a reviewer can reproduce the decision and see exactly which finding decided it.

\subsection*{The pipeline is reproducible offline}

All \todo{55} frozen \rfive decisions were replayed from the frozen upstream, case-fit and clinical inputs and were byte-identical to the delivered files. The 41-patient direct-prompting comparison was likewise rescored from \todo{164} frozen model outputs with identical metrics. The public code passes \todo{319} Python and \todo{8} Node offline tests that cover format processing, the deterministic scorer, the review, literature, case-fit and clinical rules, \rfive, rationale handoff, delivery and multi-cohort patient-identity isolation; the test harness blocks network access and strips API keys, so no test depends on a paid model or on live PubMed. The current scorer (\vtwenty) reproduces the picked set of every one of the \todo{55} patients scored with the version used for the reported results (v19), while changing candidate levels in \todo{15} and pick order in \todo{5} patients (Supplementary Note~\todo{4}).
